% =====================================================================
\chapter{Background Theory and Literature Review}
\label{ch:background}
% =====================================================================

This chapter establishes the foundations on which the thesis builds. It proceeds
from the application domain (document understanding and information extraction),
through the continual-learning problem and its method families, to the specific
and still-nascent literature on continual learning for information extraction. It
then reviews what is known about \emph{where} forgetting occurs in multimodal
models, surveys the backbones relevant to continual document IE, and closes by
isolating the research gaps that motivate the contributions of
Chapter~\ref{ch:methodology}.

\section{Document Understanding and Information Extraction}
\label{sec:bg-document-ie}

\subsection{Information-Extraction Tasks}
\label{sec:bg-ie-tasks}

Information extraction (IE) converts unstructured content into structured records.
Its canonical sub-tasks are named-entity recognition (NER), which labels spans
with semantic types; relation extraction (RE), which links entity pairs; and
event or key-information extraction (KIE), which populates schema fields. In the
document setting these tasks differ from their sentence-level counterparts in
three ways that matter for continual learning. First, documents are
\emph{spatially structured}: the role of a token (a form's question versus its
answer, a receipt's label versus its total) is often signalled by position and
typography rather than by surrounding words. Second, documents are
\emph{multimodal}: text, layout, and image jointly determine meaning. Third,
document schemas are \emph{diverse and evolving}: a forms dataset and a receipts
dataset share little vocabulary and almost no label space, so a system that
learns them sequentially faces large distribution and output-space shifts.

This thesis adopts the dominant framing of KIE as token-level sequence labelling
with the BIO tagging scheme, where each OCR token is assigned a Begin/Inside/Outside
tag for its entity type and the model is trained with token-level cross-entropy.
This framing aligns with the mainstream visually-rich document understanding
literature, yields a unified output format across datasets, and is directly
compatible with the encoder-only backbones used here.

\subsection{Multimodal Document Encoders}
\label{sec:bg-multimodal-encoders}

A multimodal document encoder ingests three streams---OCR tokens, their bounding
boxes (2D position embeddings), and image patches---and fuses them through
transformer self-attention. LayoutLMv3~\citep{huang2022layoutlmv3} unifies text
and image masking with a word--patch alignment objective and is the
field's reference encoder-only backbone. LiLT~\citep{wang2022lilt} decouples
layout from text so that a language-agnostic layout transformer can be paired with
different text encoders. BROS~\citep{hong2022bros} emphasises relative spatial
encoding and an area-masking objective. DocFormer~\citep{appalaraju2021docformer}
couples modalities through a shared spatial attention. A key structural feature
shared by these models, and central to this thesis, is that they contain
\emph{distinguishable components}---a text stream, a visual stream, layout
embeddings, and cross-modal fusion---whose individual behaviour under continual
learning has not previously been measured. Figure~\ref{fig:layoutlmv3-components}
annotates these four components on the LayoutLMv3 architecture.

\begin{figure}[t]
  \centering
  % TODO: add figure file (e.g. figures/layoutlmv3_components.pdf)
  % \includegraphics[width=0.9\linewidth]{layoutlmv3_components.pdf}
  \fbox{\parbox[c][4cm][c]{0.9\linewidth}{\centering\itshape
  Placeholder --- LayoutLMv3 architecture annotated with the four diagnosed
  components: text stream, visual/patch stream, layout (2D position) embeddings,
  and cross-modal fusion layers.}}
  \caption[LayoutLMv3 tri-modal architecture and its four components]{The
  LayoutLMv3 tri-modal architecture. OCR tokens, 2D bounding-box position
  embeddings, and image patches are projected into a shared sequence and fused by
  transformer self-attention. The four components diagnosed in this thesis---text
  stream, visual stream, layout embeddings, and cross-modal fusion---are
  highlighted; their individual susceptibility to forgetting is the object of
  Chapter~\ref{ch:methodology}.}
  \label{fig:layoutlmv3-components}
\end{figure}

\section{Catastrophic Forgetting and Continual Learning}
\label{sec:bg-cl}

\subsection{The Stability--Plasticity Dilemma}
\label{sec:bg-stability-plasticity}

Continual learning is governed by the \emph{stability--plasticity dilemma}: a
model must remain \emph{plastic} enough to acquire new tasks yet \emph{stable}
enough to retain old ones. Pure fine-tuning maximises plasticity and suffers
catastrophic forgetting; freezing the network maximises stability but prevents
learning. Every continual-learning method is, in effect, a particular point on
this trade-off, and a recurring theme of this thesis is that the right operating
point may differ \emph{per component} of a multimodal encoder.

\subsection{Continual-Learning Scenarios}
\label{sec:bg-cl-scenarios}

The literature distinguishes three canonical scenarios by what changes across the
task sequence and what is known at test time. In \emph{task-incremental} learning
(TIL) the task identity is available at inference, so a separate head can be
selected. In \emph{domain-incremental} learning (DIL) the input distribution
shifts but the label space is shared and task identity is unknown. In
\emph{class-incremental} learning (CIL) new classes are introduced over time and
the model must discriminate among all classes seen so far without task identity;
CIL is generally the hardest because the output space itself grows. For document
KIE, CIL additionally provokes the \emph{semantic shift of the ``O'' tag}: tokens
that were background (O) in an earlier session may become a labelled entity in a
later one, so the meaning of ``O'' is non-stationary. This thesis studies CIL,
DIL, and a mixed scenario, and explicitly leaves TIL to future work.

\subsection{The Five Method Families}
\label{sec:bg-cl-families}

Following the widely-used taxonomy of~\citet{wang2024survey}, continual-learning
methods fall into five families summarised in Table~\ref{tab:cl-families}.
\emph{Regularisation-based} methods (EWC~\citep{kirkpatrick2017ewc},
SI, MAS) penalise changes to parameters deemed important for past tasks, typically
weighted by a Fisher-information estimate. \emph{Replay-based} methods
(GEM~\citep{lopezpaz2017gem}, ER~\citep{rolnick2019er},
DER++~\citep{buzzega2020der}) store and re-present a memory of past exemplars.
\emph{Optimisation-based} methods constrain or project gradients so that updates
do not interfere with past tasks. \emph{Representation-based} methods freeze a
large pre-trained backbone and learn small task-specific components---prompts
(L2P~\citep{wang2022l2p}, DualPrompt~\citep{wang2022dualprompt},
CODA-Prompt~\citep{smith2023coda}) or low-rank adapters
(O-LoRA~\citep{wang2023olora}). \emph{Architecture-based} methods allocate
separate parameters per task. The dominant paradigm in 2024--2026 is
representation-based continual learning with pre-trained models: freeze the
backbone and learn lightweight, often orthogonal, task components.

\begin{table}[t]
  \centering
  \caption[The five continual-learning method families]{The five
  continual-learning method families~\citep{wang2024survey}, their core
  mechanism, and representative methods. The methods in bold are used as
  baselines in this thesis (Chapter~\ref{ch:experimental-setup}).}
  \label{tab:cl-families}
  \begin{tabular}{lll}
    \toprule
    \textbf{Family} & \textbf{Core mechanism} & \textbf{Representative methods} \\
    \midrule
    Regularisation   & Parameter-importance penalties & \textbf{EWC}, \textbf{LwF}, SI, MAS \\
    Replay           & Store and replay past exemplars & \textbf{GEM}, \textbf{ER}, \textbf{DER++}, iCaRL \\
    Optimisation     & Gradient projection/constraints & GPM, OGD, A-GEM \\
    Representation    & Frozen PTM + prompts/adapters & \textbf{L2P}, \textbf{DualPrompt}, \textbf{CODA-Prompt}, \textbf{O-LoRA} \\
    Architecture     & Separate parameters per task & PackNet, Progressive Nets \\
    \bottomrule
  \end{tabular}
\end{table}

\section{Continual Learning for Information Extraction}
\label{sec:bg-cl-ie}

Continual learning for information extraction is real but nascent, fragmented, and
almost entirely sentence-level. Table~\ref{tab:cl-ie-landscape} maps the landscape.

\subsection{Continual NER and the O-Label Problem}
\label{sec:bg-continual-ner}

CPFD~\citep{zhang2023cpfd} addresses continual NER with pooled-feature
distillation and pseudo-labelling of old entities, reporting a clear improvement
over prior continual-NER methods but remaining class-incremental on a fixed
schema. GenCNER reformulates NER \emph{generatively}, sidestepping the O-tag
semantic-shift problem by generating entity spans rather than classifying every
token; however, it uses a small text-only backbone and no document context. The
shared limitation is that both operate at sentence level and ignore layout and
vision---precisely the modalities that distinguish document IE.

\subsection{Continual Relation Extraction}
\label{sec:bg-continual-re}

WAVE++ tackles continual relation extraction with task-specific prompt pools in a
mixture-of-experts style under a task-incremental setting on TACRED/FewRel.
FCRE-OIE extends to open-world relations through open information extraction and
knowledge-graph completion in a few-shot regime. Both remain closed-world or
sentence-level and do not address document-level, cross-sentence relations.

\subsection{Document-Level Continual Learning}
\label{sec:bg-doc-cl}

Document-level continual learning is the least explored and most relevant to this
thesis. GDI-Bench together with the LW-AFT method provides the first document
intelligence continual-learning benchmark and shows that supervised fine-tuning of
a large multimodal model (InternVL3-8B) causes measurable forgetting, especially
on harder reasoning tasks; LW-AFT mitigates this by \emph{heuristically} freezing
layers. Separately, continual table detection demonstrates continual learning for
a document detection task using simple experience replay. The crucial gap is that
these works either use a heuristic remedy or address detection rather than IE, and
none diagnoses \emph{which component} of the backbone forgets---the question this
thesis answers.

\begin{table}[t]
  \centering
  \caption[The continual-learning-for-IE landscape]{The continual-learning-for-IE
  landscape and the gap each line leaves open. Sentence-level, text-only, and
  heuristic approaches dominate; no prior work provides a per-component forgetting
  diagnosis for a tri-modal document encoder.}
  \label{tab:cl-ie-landscape}
  \begin{tabular}{lll}
    \toprule
    \textbf{Setting} & \textbf{Method} & \textbf{Gap remaining} \\
    \midrule
    Continual NER     & CPFD              & Class-IL only; fixed schema; text-only \\
    Continual NER     & GenCNER           & Small backbone; no document/layout context \\
    Continual RE      & WAVE++            & Task-IL; closed-world; sentence-level \\
    Continual RE      & FCRE-OIE          & Few-shot; no document-level relations \\
    Document CL       & GDI-Bench + LW-AFT & Heuristic freezing; no IE tasks; no diagnosis \\
    Document CL       & CL Table Detection & Detection only; simple replay; no IE \\
    Multimodal CL     & CL-CrossVQA/CLiMB & VQA/VL only; no document IE \\
    Multimodal CL     & SEFE              & Instruction tuning; not IE; no per-component view \\
    \bottomrule
  \end{tabular}
\end{table}

\section{Multimodal and Component-Level Forgetting}
\label{sec:bg-multimodal-forgetting}

\subsection{Where Forgetting Lives: Text vs.\ Visual vs.\ Fusion}
\label{sec:bg-where-forgetting}

Evidence from multimodal continual-learning benchmarks suggests that forgetting is
not uniform across modalities, but the picture is incomplete. CL-CrossVQA finds
that dual-stream visual-question-answering models forget less than single-stream
ones and that replay is the most robust remedy, hinting that architectural
structure shapes forgetting. CLiMB observes that continual learning can prevent
forgetting on vision-and-language tasks yet fails to enable positive transfer.
These findings concern VQA rather than document IE and do not localise forgetting
to specific \emph{components} of a single encoder. The methodological inspiration
for this thesis comes instead from analyses of forgetting in unimodal networks:
the ``anatomy of catastrophic forgetting''~\citep{ramasesh2021anatomy} shows that
forgetting in deep networks concentrates in higher layers and can be probed
layer-by-layer, and the study of forgetting in multimodal large language
models~\citep{zhai2023investigating} adapts such diagnostics to multimodal
settings. This thesis transfers that per-layer, per-component diagnostic stance to
a tri-modal document encoder for the first time.

\subsection{Superficial versus Essential Forgetting}
\label{sec:bg-superficial-essential}

SEFE introduces a distinction between \emph{superficial} forgetting---loss of the
expected answer format---and \emph{essential} forgetting---loss of the underlying
knowledge. This distinction matters for KIE, where a model may retain the ability
to localise entities yet lose the correct label mapping (or vice versa). It
motivates measuring forgetting with multiple, complementary instruments rather
than a single accuracy number, a principle adopted in
Chapter~\ref{ch:methodology}.

\subsection{Measuring Representational Drift: CKA}
\label{sec:bg-cka}

To attribute forgetting to components, this thesis measures how internal
representations change between checkpoints. Centered Kernel Alignment
(CKA)~\citep{kornblith2019similarity} compares two activation matrices and returns
a similarity in $[0,1]$ that is invariant to orthogonal transforms and isotropic
scaling, properties that make it more reliable than naive cosine similarity
between weights. A CKA of one indicates unchanged representations; a low CKA
indicates reorganisation. Per-layer CKA between the model before and after a task
thus localises the ``forgetting hotspots,'' and is one of the two primary
diagnostic metrics used here, alongside per-group Fisher information.

\section{Backbones for Continual Document IE}
\label{sec:bg-backbones}

The choice of backbone determines both the cost and the interpretability of the
diagnosis. LayoutLMv3-base ($\approx$\,133M parameters, encoder-only) is the
primary backbone: it is document-specific, small enough to run a full grid of
methods $\times$ scenarios $\times$ seeds on a single GPU, and exposes
clearly-separable text, visual, layout, and fusion components. Large multimodal
LLMs such as InternVL3-8B are used in document continual-learning benchmarks
(GDI-Bench) but their scale and decoder-only design make a controlled
per-component study expensive and harder to interpret; they are noted as a
comparison target rather than a primary subject. To test that the diagnosis
generalises beyond a single architecture, LiLT-base and BROS-base are used as
secondary backbones on a reduced method subset.

\section{Research Gaps and Opportunities}
\label{sec:bg-gaps}

The survey above exposes the gaps the thesis addresses:

\begin{enumerate}
  \item \textbf{No per-component forgetting diagnosis for tri-modal document
        encoders.} Prior work measures forgetting at the level of whole tasks or,
        at most, layers of unimodal models. No study attributes forgetting to the
        text, visual, layout, and fusion components of a document encoder.
  \item \textbf{Document-level versus sentence-level continual IE.} CPFD, GenCNER,
        and WAVE++ are sentence-level and text-only; document IE requires layout
        and vision and exhibits the O-tag semantic shift at document scale.
  \item \textbf{Heuristic versus principled continual learning for document
        backbones.} The only document-level remedy, LW-AFT, freezes layers
        heuristically. Whether a freezing/regularisation mechanism \emph{targeted
        by measurement} outperforms uniform or heuristic treatment is open.
  \item \textbf{Fragmented evaluation.} Existing results use different benchmarks
        and protocols; there is no shared continual benchmark for document IE that
        compares image-CL baselines and IE-specific methods on equal footing.
\end{enumerate}

\section{Positioning of the Current Research}
\label{sec:bg-positioning}

This thesis is a \emph{diagnostic-plus-remedy} study in the spirit
of~\citet{ramasesh2021anatomy} and~\citet{zhai2023investigating}: it first
characterises how forgetting distributes across the components of a tri-modal
document encoder, then proposes a method whose mechanism is concentrated on the
dominant-forgetting component. The measurement contribution stands on its own---it
is informative even if forgetting turns out to be uniform---while the remedy is
explicitly derived from, and validated against, the measurement. The work is
deliberately scoped to encoder-only document KIE on three public datasets, with
larger backbones, additional scenarios, and multilingual extensions left as future
directions (Chapter~\ref{ch:conclusions}).

\section{Chapter Summary}
\label{sec:bg-summary}

Document IE is a multimodal, spatially-structured, schema-diverse problem solved
by tri-modal encoders such as LayoutLMv3. Continual learning of such tasks is
obstructed by catastrophic forgetting, for which five method families offer
remedies---but almost all treat the network uniformly and were validated on
unimodal benchmarks. Continual learning for IE exists only in fragmented,
sentence-level, mostly text-only form, and the one document-level remedy is
heuristic. Crucially, no prior work measures \emph{which component} of a
multimodal document encoder forgets. The next chapter develops a diagnostic to
answer that question and a mechanism-targeted remedy derived from the answer.
